\section{First-order resolution}
\label{sec:fol-cardinality-normalforms-04}

\begin{definition}[First-order clause and formula interpretation]
Using the first-order literal defined before the prenex theorem, a
\emph{clause} is a finite set of first-order literals, with
set equality as its equality relation.  Thus literal order and multiplicity
are not clause data.

For an exact map back to formulas, take as presentation data a well-order
\(\prec_{\Lang}\) of the symbol tokens of \(\Lang\).  Supplying this order is
an explicit well-ordering assumption; for an arbitrary signature,
its existence follows from the well-ordering theorem and hence from the Axiom
of Choice.  The semantic results below do not depend on which supplied
well-order is used.  Combine \(\prec_{\Lang}\) with the order
\(x_0,x_1,\ldots\) and a fixed order of
the finitely many logical and punctuation tokens, and order formula strings
first by length and then lexicographically.  For a clause \(C\), let
\(k=\abs{C}\); when \(k>0\), list its members as
\(L_1\prec\cdots\prec L_k\) in the induced order and define, by downward
recursion,
\[
 D_k\defeq L_k,
 \qquad
 D_i\defeq(L_i\lor D_{i+1})\quad(1\leq i<k).
\]
Then define
\[
 \operatorname{form}_{\Lang}(C)\defeq
 \begin{cases}
  \bottom,&C=\varnothing,\\
  D_1,&C\ne\varnothing.
 \end{cases}
\]
Let \(x_{i_1},\ldots,x_{i_m}\), with
\(i_1<\cdots<i_m\), be exactly the variables occurring in \(C\).  Its
formula interpretation is
\[
 \operatorname{ucl}_{\Lang}(C)\defeq
 \forall x_{i_1}\cdots\forall x_{i_m}\,
 \operatorname{form}_{\Lang}(C),
\]
with the quantifier prefix empty when \(m=0\).  The functions
\(\operatorname{form}_{\Lang}\) and \(\operatorname{ucl}_{\Lang}\) are
metalinguistic, not new object-language constructors.  They are required
because a clause is a set rather than a formula;
\(\operatorname{ucl}_{\Lang}\) differs from \(\operatorname{form}_{\Lang}\)
by closing every clause variable.  A structure satisfies \(C\) iff it
satisfies \(\operatorname{ucl}_{\Lang}(C)\), and it satisfies a clause set
iff it satisfies every member.

A clause is \emph{equality-free} iff no literal in it contains \(=\); that is,
the atom underlying each of its literals has the form
\(R(t_1,\ldots,t_n)\), not \(r=t\).  A clause set is equality-free iff each
of its clauses is equality-free.  This condition concerns the literals in
the clauses, not the signature, because logical equality is not a signature
member.
\end{definition}

\begin{definition}[Finite-support simultaneous term substitution]
A \emph{term substitution} is a total function
\(\theta:\Var\to\Term(\Lang)\) whose support
\[
 \operatorname{supp}(\theta)
 \defeq\set{x\in\Var:\theta(x)\ne x}
\]
is finite.  A finite display such as
\(\set{x_1\mapsto t_1,\ldots,x_n\mapsto t_n}\) denotes the substitution
shown on those variables and the identity on every other variable.

Its simultaneous homomorphic action is defined by
\[
 x\theta=\theta(x),
 \qquad
 f(t_1,\ldots,t_n)\theta
 =f(t_1\theta,\ldots,t_n\theta),
\]
\[
 \begin{aligned}
 R(t_1,\ldots,t_n)\theta
   &\defeq R(t_1\theta,\ldots,t_n\theta),\\
 (r=t)\theta&\defeq(r\theta=t\theta),\\
 (\neg A)\theta&\defeq\neg(A\theta),\\
 C\theta&\defeq\set{L\theta:L\in C}.
 \end{aligned}
\]
Here \(A\) is atomic, \(L\) is a literal, and \(C\) is a clause.  The
application is simultaneous: when \(\theta(x)\) is inserted for \(x\), the
same application does not apply \(\theta\) again inside that inserted term.
The clause action may identify two literals, in which case set equality
removes the duplicate.  This quantifier-free action is the one needed by
resolution; it neither replaces nor restricts the total capture-avoiding
single-variable formula substitution \(\varphi[t/x]\) defined earlier.
\end{definition}

\begin{definition}[Composition and instance preorder]
The identity substitution \(\mathrm{id}\) maps every variable to itself.  For
substitutions \(\theta\) and \(\eta\), their composition \(\theta\eta\) is
oriented by
\[
 (\theta\eta)(x)\defeq(\theta(x))\eta.
\]
Its support is contained in
\(\operatorname{supp}(\theta)\cup\operatorname{supp}(\eta)\), so it is
again a finite-support term substitution.
Thus \(t(\theta\eta)=(t\theta)\eta\), and the same equation holds for
atoms, literals, and clauses.  Define the \emph{instance preorder} on term
substitutions by
\[
 \theta\preceq_{\mathrm{inst}}\rho
 \quad\Longleftrightarrow\quad
 \text{there is a term substitution \(\eta\) such that }
 \rho=\theta\eta.
\]
This is a relation among substitutions, not equality: it records that
\(\rho\) results from further simultaneous instantiation of \(\theta\).
It is reflexive by \(\mathrm{id}\) and transitive by associativity of the
displayed composition, so it is a preorder.  Its use below is precisely to
state what ``most general'' means.
\end{definition}

\begin{lemma}[Composition of simultaneous substitutions]
For every term \(t\), atom \(A\), literal \(L\), and clause \(C\),
\[
 \begin{gathered}
 t(\theta\eta)=(t\theta)\eta,
 \qquad A(\theta\eta)=(A\theta)\eta,\\
 L(\theta\eta)=(L\theta)\eta,
 \qquad C(\theta\eta)=(C\theta)\eta.
 \end{gathered}
\]
\end{lemma}
\begin{proof}
Induct on terms for the first equation.  The atomic and literal equations then
follow from their displayed homomorphic clauses.  Taking the pointwise image
of a finite set proves the clause equation.  The term equation also proves
associativity of substitution composition by extensionality.
\end{proof}

\begin{lemma}[Satisfaction under simultaneous term substitution]
\label{lem:clause-substitution}
If a \(\Lang\)-structure \(\Meta\) satisfies a clause \(C\), then it
satisfies \(C\theta\) for every term substitution \(\theta\).
\end{lemma}
\begin{proof}
For an assignment \(s\), define another assignment by
\[
 s^{\theta}(x)\defeq\sem{\theta(x)}^{\Meta}_{s}.
\]
Induction on terms gives
\(\sem{t\theta}^{\Meta}_{s}=\sem{t}^{\Meta}_{s^{\theta}}\), and the atomic
truth clauses extend this equality to literals.  Because \(\Meta\) satisfies
the universal closure of \(C\), the assignment \(s^{\theta}\) satisfies some
literal \(L\in C\).  Hence \(s\) satisfies \(L\theta\in C\theta\).  Since
\(s\) was arbitrary, \(\Meta\) satisfies the universal closure of
\(C\theta\).
\end{proof}

\begin{definition}[Renaming apart]
A \emph{variable renaming} is a finite-support term substitution induced by a
bijection \(\pi:\Var\to\Var\); hence every \(\pi(x)\) is a variable.  For a
clause \(C\), put
\[
 \operatorname{Var}(C)\defeq
 \bigcup_{L\in C}\operatorname{Var}(L).
\]
Clauses \(C\pi\) and \(D\rho\) are \emph{renamed apart} iff
\(\operatorname{Var}(C\pi)\cap\operatorname{Var}(D\rho)=\varnothing\).
Before a resolution inference, each parent is replaced by such a renamed
copy.  The required renamings exist because both variable sets are finite and
\(\Var\) is infinite.  Universal closure makes a renamed copy equivalent to
the original clause; the operation changes bound variable names in its
formula interpretation, not its satisfaction conditions.
\end{definition}

\begin{definition}[Unifier and most general unifier]
A term substitution \(\theta\) is a \emph{unifier} of atomic formulas
\(A\) and \(B\) iff \(A\theta\) and \(B\theta\) are literally identical.
A unifier \(\theta\) is a \emph{most general unifier} iff
\(\theta\preceq_{\mathrm{inst}}\rho\) for every unifier \(\rho\); equivalently, every
unifier \(\rho\) has the form \(\theta\eta\) for some term substitution
\(\eta\).
\end{definition}

\begin{definition}[First-order unification procedure]
Below, \emph{atom} abbreviates \emph{atomic formula}.  To unify atoms
\(A\) and \(B\), first fail if their outer atomic symbols or
arities differ.  Otherwise form a finite set \(E\) of metalinguistic term
equations \(r\doteq t\) by pairing corresponding arguments (the two ordered
arguments when the common atomic symbol is \(=\)), and initialize
\(\theta\) to \(\mathrm{id}\).  Repeatedly apply the following transformations
to a selected equation:
\begin{enumerate}
  \item delete \(t\doteq t\);
  \item replace
        \(f(r_1,\ldots,r_n)\doteq f(t_1,\ldots,t_n)\) by the equations
        \(r_i\doteq t_i\) for \(1\leq i\leq n\);
  \item fail on
        \(f(r_1,\ldots,r_m)\doteq g(t_1,\ldots,t_n)\) when
        \(f\ne g\) or \(m\ne n\);
  \item orient \(t\doteq x\) as \(x\doteq t\) when \(x\) is a variable
        and \(t\) is not;
  \item for \(x\doteq t\) with \(x\ne t\), perform the \emph{occurs
        check}: if \(x\in\operatorname{Var}(t)\), fail; otherwise put
        \(\delta=\set{x\mapsto t}\), apply \(\delta\) simultaneously to
        every remaining equation, and replace the accumulator by
        \(\theta\delta\).
\end{enumerate}
If the equation set becomes empty, return \(\theta\).  The occurs check is
part of this procedure rather than part of the definition of a term
substitution: it rejects the impossible demand that a finite term equal a
proper term containing itself.
\end{definition}

\begin{theorem}[First-order unification]
The unification procedure terminates.  It returns failure iff its
input atoms are not unifiable; otherwise it returns a most general unifier.
\end{theorem}
\begin{proof}[Proof sketch]
Maintain the invariant that the unifiers of the original problem are exactly
the substitutions \(\theta\eta\) for substitutions \(\eta\) solving every
equation still in \(E\).  Deletion, decomposition, and orientation preserve
this invariant directly.  Elimination preserves it by the composition lemma.
The clash rule has no solution.  If \(x\) occurs properly in \(t\), no finite
term can satisfy \(x\eta=t\eta\), which justifies occurs-check failure.

An elimination removes its variable from the remaining problem; a
decomposition removes matching outer function symbols; deletion removes an
equation, and orientation is immediately followed by one of these steps.
The corresponding lexicographic measure therefore decreases, proving
termination.  At success the remaining equation set is empty, so the
invariant says both that \(\theta\) unifies the original atoms and that every
other unifier is an instance \(\theta\eta\).  At either failure rule the
invariant says that no unifier exists.  The full lexicographic-measure check
and the transformation-by-transformation invariant calculations are omitted.
\end{proof}

After the two parent clauses \(E\) and \(F\) have been renamed apart, choose
\(A\in E\) and \(\neg B\in F\), and let \(\theta\) be a most general
unifier of the atoms \(A\) and \(B\).  The first-order resolution rule is
\[
 \Rule[Resolution]{E\qquad F}
 {\bigl((E\setminus\set{A})\cup(F\setminus\set{\neg B})\bigr)\theta}.
\]
Factoring applies to two distinct literals of the same polarity in one
clause \(E\).  If their underlying atoms are \(A\) and \(B\), and
\(\theta\) is a most general unifier of those atoms, its rule is
\[
 \Rule[Factoring]{E}{E\theta}.
\]
Because clauses are sets, the two selected literals become one member of
\(E\theta\).  A \emph{first-order resolution derivation} from a clause set
\(S\) is a finite sequence in which each clause belongs to \(S\) or follows
by resolution or factoring from earlier clauses, using renamed-apart copies
of two parents before every resolution step.

\begin{theorem}[Soundness of first-order resolution and factoring]
Every structure satisfying the parent clauses of a resolution or factoring
inference satisfies its conclusion.  Consequently, if a first-order
resolution derivation from \(S\) contains the empty clause, then \(S\) is not
satisfiable.
\end{theorem}
\begin{proof}
Renaming apart preserves the universally closed formula interpretations: one
direction follows from \cref{lem:clause-substitution}, and the other follows
by applying the inverse variable renaming.  By the same lemma, a structure
satisfying \(E\) and \(F\) satisfies \(E\theta\) and \(F\theta\).  If an
assignment falsified every literal in the displayed resolvent, satisfaction
of the substituted first parent would force \(A\theta\), while satisfaction
of the substituted second parent would force \(\neg(B\theta)\).  But
\(A\theta=B\theta\) literally, a contradiction.  Factoring is sound directly
by \cref{lem:clause-substitution}.  Induction on a derivation now preserves
satisfaction, whereas the formula interpretation of the empty clause is
\(\bottom\).
\end{proof}

\begin{definition}[Ground term and ground instance]
A term \(t\) is \emph{ground} iff \(\operatorname{Var}(t)=\varnothing\).  Let
\(C\) be a first-order clause.  A \emph{ground instance} of \(C\) is a clause
\(C\theta\) where \(\theta\) maps every variable occurring in \(C\) to a
ground term.  These terms classify terms and substitution instances,
respectively; neither means a closed formula.  They are required in the two
cited results and the resolution-completeness proof below.
\end{definition}

\begin{theorem}[Herbrand theorem for clauses]
\label{thm:herbrand-clauses}
Let \(S\) be a finite equality-free set of first-order clauses in a signature
having at least one ground term.  Then \(S\) is not satisfiable iff some finite
set of ground instances of members of \(S\), regarded as propositional
clauses whose distinct ground atoms are propositional variables, is not
satisfiable.
\textit{Status: cited theorem.}
\end{theorem}

\begin{theorem}[Lifting lemma]
\label{thm:resolution-lifting}
Let \(C,D\) be first-order clauses with disjoint variable sets, and let \(G\)
be the ground resolvent produced by one ground resolution step from ground
instances of \(C\) and \(D\).  Factor each parent literal family that its
ground substitution identifies, and let \(R\) be the first-order resolvent of
the resulting parents obtained with a most general unifier.  Then there is a
substitution \(\eta\) that maps every variable occurring in \(R\) to a ground
term and satisfies \(R\eta=G\).  The analogous equality holds for a ground
factoring step and its lifted first-order factoring step.  Consequently, for
every clause set \(S\), every
 finite ground resolution refutation of ground instances of \(S\) lifts, step by
step, to a first-order resolution derivation with factoring from \(S\).
\textit{Status: cited theorem.}
\end{theorem}

See \cite{Robinson1965,BachmairGanzinger2001} for complete proofs of
\cref{thm:herbrand-clauses,thm:resolution-lifting}.

\begin{example}[Unification and resolution]
Resolve
\[
 \set{P(f(x)),Q(x)}
 \quad\text{and}\quad
 \set{\neg P(y),R(y)}.
\]
The substitution \(\theta=\set{y\mapsto f(x)}\) is a most general unifier of
\(P(f(x))\) and \(P(y)\).  The resolvent is
\[
 \set{Q(x),R(f(x))}.
\]
\end{example}

\begin{theorem}[Refutational completeness of first-order resolution]
\label{thm:fol-resolution}
Let \(S\) be a finite equality-free set of first-order clauses.  After adding
a constant if the signature has no ground term, \(S\) is not satisfiable iff
resolution with factoring derives the empty clause.
\end{theorem}
\begin{proof}[Proof sketch]
The preceding soundness theorem proves the reverse implication.  For the
forward implication, \cref{thm:herbrand-clauses} yields a finite propositionally
unsatisfiable set of ground instances of clauses in \(S\).  Propositional
refutational completeness gives a ground resolution refutation.
\Cref{thm:resolution-lifting} replaces each ground resolution step by a first-order resolution or
factoring step whose instance is that ground step.  Induction through the
finite ground refutation yields a first-order derivation of the empty clause.
This proof depends on the two cited theorems, the existence of most general
unifiers for unifiable first-order terms, and propositional refutational
completeness; it is the
precise equality-free scope of the classical result in
\cite{Robinson1965,BachmairGanzinger2001}.  The cited complete proofs and the
step-by-step lifted derivation are not reproduced here.
\end{proof}

Bare resolution is not complete for logical equality.  One must add
an equality-complete treatment, such as equality axioms together with
resolution or a sound and complete paramodulation or superposition calculus.
The theorem above deliberately excludes every clause having a literal that
contains logical equality.

